\documentclass[11pt]{amsart}
\usepackage{amsmath,amssymb,amsthm}
\usepackage[margin=1.1in]{geometry}
\usepackage{array,longtable,booktabs}
\usepackage{hyperref}

\newtheorem{theorem}{Theorem}
\newtheorem{lemma}[theorem]{Lemma}
\newtheorem{corollary}[theorem]{Corollary}
\newtheorem{proposition}[theorem]{Proposition}
\theoremstyle{remark}
\newtheorem{remark}[theorem]{Remark}

\newcommand{\Inf}{\mathrm{Inf}}
\newcommand{\E}{\mathbb{E}}

\title{The one-coordinate step in the Fourier entropy-influence conjecture}
\author{Yitzchak Shmalo}
\address{Einstein Institute of Mathematics, The Hebrew University of Jerusalem, Jerusalem, Israel}
\date{}

\begin{document}

\begin{abstract}
For a Boolean function $f$ and a coordinate $i$ let $D_i(f)=H(f)-\tfrac12 H(f_{i\to 1})-\tfrac12 H(f_{i\to -1})$ be the
drop in Fourier entropy when $x_i$ is fixed. Gonz\'alez, MacManus and Pereyra observed that the Fourier
entropy-influence conjecture follows, by induction on the dimension, if every Boolean function has a coordinate
with $D_i(f)\le K\,\Inf_i(f)$. We show that this condition is also necessary: if the conjecture holds with constant
$C$, then every nonconstant Boolean function has a coordinate whose erasure entropy is at most $C\,\Inf_i(f)$, and
hence one with $D_i(f)\le (C+2/\ln 2)\Inf_i(f)$. The constant $2/\ln 2$ is the sharp bound for the Jensen--Shannon
part of $D_i$. We also give an equivalent pseudo-likelihood form of the conjecture, a bound on the influences of
extremal functions, and an example showing that the condition cannot be required at every coordinate.
\end{abstract}

\subjclass[2020]{06E30, 42C10, 94A17}
\keywords{Fourier entropy-influence conjecture, Boolean functions, Fourier entropy, total influence}

\maketitle

\section{Introduction}\label{sec:intro}

Let $f:\{-1,1\}^n\to\{-1,1\}$ have Fourier expansion $f=\sum_S \hat f(S)\chi_S$. By Parseval, $\mu(S)=\hat f(S)^2$
is a probability distribution on subsets of $[n]$, the spectral distribution of $f$. Write
\[
H(f)=\sum_S \mu(S)\log_2\frac{1}{\mu(S)},\qquad \Inf_i(f)=\sum_{S\ni i}\mu(S),\qquad I(f)=\sum_i \Inf_i(f)=\sum_S|S|\mu(S).
\]
The Fourier entropy-influence (FEI) conjecture of Friedgut and Kalai~\cite{FK96} asks for a universal $C$ with
$H(f)\le C\,I(f)$ for every Boolean $f$. It is known for symmetric functions and read-once decision
trees~\cite{OWZ11}, for read-once formulas~\cite{OT13,CKLS} and for functions of bounded average decision-tree
depth~\cite{WWW14}. The conjecture was proved for linear threshold functions in~\cite{ShmLTF} and for
plateaued functions in~\cite{Pot26}. For every
Boolean $f$, Han~\cite{Han} proved $H(f)\le(2+3/\ln 2)\,I(f)+\sum_i\Inf_i(f)\log_2(1/\Inf_i(f))$, and Li and
Han~\cite{LH25} lowered the constant in front of $I(f)$ to the sharp value $1/\ln 2$, with a proof that holds for
every real-valued $f$ with $\|f\|_2=1$. The min-entropy version is studied in~\cite{ACKSW,KKLMS}.

Fix a coordinate $i$ and the restrictions $f^{\pm}=f_{i\to\pm1}$. Since $I(f)=\tfrac12 I(f^+)+\tfrac12 I(f^-)+\Inf_i(f)$,
the inequality
\begin{equation}\label{eq:step}
D_i(f):=H(f)-\tfrac12 H(f^+)-\tfrac12 H(f^-)\le K\,\Inf_i(f)
\end{equation}
for one coordinate $i$ of every Boolean $f$ gives $H(f)\le K\,I(f)$ by induction on $n$. This route is suggested by
Han~\cite[(3.17) and the discussion after Remark 3.2]{Han} in a different local form, and is set out by Gonz\'alez, MacManus and Pereyra~\cite{GMP},
who prove~\eqref{eq:step} along suitable decision trees for several new classes. The inequality version
$D_1(f)\le 2$ is Lemma~3.2 of Chakraborty, Kulkarni, Lokam and Saurabh~\cite{CKLS}.

Our main observation is that the route loses nothing: the one-coordinate condition is equivalent to the conjecture.
Let $C_n=\max H(f)/I(f)$ over nonconstant Boolean $f$ on $n$ variables, and
\[
R_n=\max_f \min_{i:\,\Inf_i(f)>0}\frac{E_i(f)}{\Inf_i(f)},\qquad
R'_n=\max_f \min_{i:\,\Inf_i(f)>0}\frac{D_i(f)}{\Inf_i(f)},
\]
with the maxima over the same functions and $E_i$ the erasure entropy defined in Section~\ref{sec:split}.

\begin{theorem}\label{thm:main}
For every $n$,
\[
R_n\le C_n\le \max_{m\le n}R_m+\frac{2}{\ln 2},\qquad
\max_{m\le n} R'_m \ge C_n,\qquad R'_n\le C_n+\frac{2}{\ln 2}.
\]
In particular the FEI conjecture holds if and only if $\sup_n R_n<\infty$, if and only if $\sup_n R'_n<\infty$.
\end{theorem}

The proof is short. The entropy drop splits into an erasure part and a Jensen--Shannon part (this is the computation
in~\cite[Appendix B]{GMP}); the Jensen--Shannon part is at most $(2/\ln2)\Inf_i$, with this constant sharp
(Lemma~\ref{lem:jsd}); and the erasure parts of all coordinates add up to at most $H(f)$
(Lemma~\ref{lem:sum}). Section~\ref{sec:more} gives the pseudo-likelihood form, the bound on extremal influences, and
the example.

\section{The split}\label{sec:split}

Let $U$ be a random set with law $\mu$ and identify $U$ with its indicator vector $(X_1,\dots,X_n)$. For
$T\subseteq[n]\setminus\{i\}$ write $a_T=\hat f(T)$ and $b_T=\hat f(T\cup\{i\})$. Then
$\widehat{f^{\pm}}(T)=a_T\pm b_T$, so $\mu^{\pm}(T)=(a_T\pm b_T)^2$ are the spectral distributions of the two
restrictions, and $\tfrac12(\mu^+(T)+\mu^-(T))=a_T^2+b_T^2=\Pr[U\setminus\{i\}=T]$. Let
$h(q)=q\log_2\frac1q+(1-q)\log_2\frac1{1-q}$. Define
\[
E_i(f)=H(X_i\mid U\setminus\{i\})=\sum_T (a_T^2+b_T^2)\,h\!\Big(\frac{b_T^2}{a_T^2+b_T^2}\Big),
\qquad
J_i(f)=\mathrm{JS}(\mu^+,\mu^-),
\]
where $\mathrm{JS}(p,q)=H(\tfrac{p+q}2)-\tfrac12H(p)-\tfrac12H(q)$ is the Jensen--Shannon divergence in bits. The quantity $E_i(f)$ is the erasure entropy of Verd\'u and
Weissman~\cite{VW08} at coordinate $i$ of the random set $U$.

\begin{lemma}[\cite{GMP}, Appendix B]\label{lem:split}
$D_i(f)=E_i(f)+J_i(f)$.
\end{lemma}

\begin{proof}
By the chain rule $H(f)=H(U)=H(U\setminus\{i\})+H(X_i\mid U\setminus\{i\})$. The law of $U\setminus\{i\}$ is
$\tfrac12(\mu^++\mu^-)$, so $H(U\setminus\{i\})=\tfrac12H(f^+)+\tfrac12H(f^-)+\mathrm{JS}(\mu^+,\mu^-)$.
\end{proof}

Both parts are nonnegative, and $E_i(f)\le H(X_i)=h(\Inf_i(f))$.

\begin{lemma}\label{lem:jsd}
$J_i(f)\le \frac{2}{\ln 2}\Inf_i(f)$ for every Boolean $f$ and every coordinate $i$, and the constant is sharp.
\end{lemma}

\begin{proof}
Write $|a_T|=r\cos\theta$, $|b_T|=r\sin\theta$ with $\theta\in[0,\pi/2]$, and $s=\sin 2\theta$. The term of $T$ in
$J_i$ is $r^2\bigl(1-h(\tfrac{1+\sigma s}2)\bigr)$ with $\sigma=\pm1$ the sign of $a_Tb_T$; since $h$ is symmetric,
it equals $r^2\varphi(s)/\ln 2$, where
\[
\varphi(s)=\tfrac12\bigl[(1+s)\ln(1+s)+(1-s)\ln(1-s)\bigr].
\]
The term of $T$ in $\Inf_i$ is $b_T^2=r^2\sin^2\theta$, so it suffices to show $\varphi(\sin2\theta)\le 2\sin^2\theta$.
For $\theta\ge\pi/4$ this holds because $\varphi\le\ln2<1\le 2\sin^2\theta$. On $[0,\pi/4]$ put
$F(\theta)=2\sin^2\theta-\varphi(\sin 2\theta)$, which is continuous. Then $F(0)=0$ and, on $[0,\pi/4)$, since
$\varphi'(s)=\operatorname{artanh}s$ and $\cos2\theta=\sqrt{1-s^2}$ there,
\[
F'(\theta)=2s-2\sqrt{1-s^2}\,\operatorname{artanh}(s).
\]
With $s=\tanh y$, $y\ge0$, this is $2\tanh y-2y/\cosh y=2(\sinh y-y)/\cosh y\ge0$. Hence $F\ge0$ on
$[0,\pi/4]$.

For sharpness, $\varphi(s)=s^2/2+O(s^4)$, so $\varphi(\sin2\theta)/(2\sin^2\theta)\to1$ as $\theta\to0$. Take a
Boolean $f$ and coordinate $i$ with $b_T^2/(a_T^2+b_T^2)\to 0$ uniformly in $T$ along a sequence, for example the
function of Proposition~\ref{prop:point} below; then $J_i/\Inf_i\to 2/\ln 2$.
\end{proof}

The bound in~\cite[Lemma 8.1]{GMP} is $1-h(t)\le 4(\tfrac12-t)^2$, which gives $J_i\le 4\,\Inf_i$.

\begin{lemma}\label{lem:sum}
$\sum_{i=1}^n E_i(f)\le H(f)$.
\end{lemma}

\begin{proof}
$H(f)=H(X_1,\dots,X_n)=\sum_i H(X_i\mid X_1,\dots,X_{i-1})\ge\sum_i H(X_i\mid X_j,\,j\ne i)=\sum_iE_i(f)$, since
conditioning on more variables does not increase entropy.
\end{proof}

Since $H(X_{[n]\setminus\{i\}})=H(f)-E_i(f)$, Lemma~\ref{lem:sum} is equivalent to Han's inequality~\cite{Han78}
\[
\textstyle\sum_iH(X_{[n]\setminus\{i\}})\ge(n-1)H(X_1,\dots,X_n).
\]

Subadditivity gives the matching upper bound $H(f)\le\sum_iH(X_i)=\sum_ih(\Inf_i(f))$; this is Claim~4.1 of Keller, Mossel and Schlank~\cite{KMS}.

\section{Proof of Theorem~\ref{thm:main}}\label{sec:proof}

\emph{$R_n\le C_n$.} Let $f$ be nonconstant on $n$ variables. By Lemma~\ref{lem:sum},
$\sum_iE_i(f)\le H(f)\le C_nI(f)=C_n\sum_i\Inf_i(f)$. The coordinates with $\Inf_i(f)=0$ have $X_i=0$ almost surely
and $E_i(f)=0$, so some coordinate with $\Inf_i(f)>0$ has $E_i(f)\le C_n\Inf_i(f)$.

\emph{$C_n\le\max_{m\le n}R_m+2/\ln2$.} Let $K=\max_{m\le n}R_m+2/\ln2$. We show $H(g)\le K\,I(g)$ for every Boolean
$g$ on $m\le n$ variables by induction on the number of coordinates $g$ depends on. If it depends on none,
$H(g)=I(g)=0$. Otherwise choose $i$ with $\Inf_i(g)>0$ and $E_i(g)\le R_m\Inf_i(g)$. By Lemmas~\ref{lem:split}
and~\ref{lem:jsd}, $D_i(g)\le K\Inf_i(g)$; the restrictions $g^{\pm}$ depend on fewer coordinates, so
\[
H(g)=\tfrac12H(g^+)+\tfrac12H(g^-)+D_i(g)\le K\bigl(\tfrac12I(g^+)+\tfrac12I(g^-)+\Inf_i(g)\bigr)=K\,I(g).
\]

\emph{$\max_{m\le n}R'_m\ge C_n$} is the same induction with $D_i$ in place of $E_i$ and no Jensen--Shannon term,
which is the argument of~\cite{GMP}. \emph{$R'_n\le C_n+2/\ln2$}: by Lemmas~\ref{lem:split}, \ref{lem:jsd}
and~\ref{lem:sum}, $\sum_iD_i(f)\le H(f)+\frac2{\ln2}I(f)\le(C_n+\frac2{\ln2})I(f)$, and a coordinate with
$\Inf_i(f)>0$ attains the average. \qed

\section{Further consequences}\label{sec:more}

\subsection*{A pseudo-likelihood form}
The quantity $\sum_iE_i(f)=\sum_i H(X_i\mid X_j,\,j\neq i)$ is the pseudo-likelihood entropy of the spectral sample:
the expected log-loss, in bits, of predicting each coordinate of $U$ from all the others, as in Besag's
pseudo-likelihood~\cite{Bes75}.

\begin{corollary}\label{cor:pl}
If $H(f)\le C\,I(f)$ for all Boolean $f$, then $\sum_iE_i(f)\le C\,I(f)$ for all Boolean $f$. Conversely, if
$\sum_iE_i(f)\le C\,I(f)$ for all Boolean $f$, then $H(f)\le(C+2/\ln 2)\,I(f)$ for all Boolean $f$.
\end{corollary}

\begin{proof}
The first statement is Lemma~\ref{lem:sum}. For the second, the hypothesis gives $R_n\le C$ for all $n$ by averaging,
and Theorem~\ref{thm:main} applies.
\end{proof}

So FEI may be read as a statement about conditional rather than joint uncertainty of the spectral sample: it is
equivalent to asking that each element of $U$ be predictable from the rest at a cost of $O(1)$ bits per expected
element.

\subsection*{Extremal functions have small influences}

\begin{proposition}\label{prop:ext}
Let $f$ be a Boolean function on $n$ variables with $H(f)=C_nI(f)$. Then every
coordinate with $\Inf_i(f)>0$ satisfies
\[
\Inf_i(f)\le 2^{\,3/\ln 2-C_{n-1}}.
\]
\end{proposition}

\begin{proof}
Let $p=\Inf_i(f)>0$. The restrictions are functions of $n-1$ variables, so $H(f^{\pm})\le C_{n-1}I(f^{\pm})$
(trivially if they are constant), and $C_{n-1}\le C_n$, so
\[
D_i(f)=C_nI(f)-\tfrac12H(f^+)-\tfrac12H(f^-)\ge C_nI(f)-C_{n-1}\bigl(I(f)-p\bigr)\ge C_{n-1}\,p .
\]
On the other side $D_i(f)\le h(p)+\frac2{\ln2}p$, and $h(p)\le p\log_2\frac1p+\frac{p}{\ln 2}$ because
$(1-p)\ln\frac1{1-p}\le p$. Dividing by $p$ gives $C_{n-1}\le\log_2\frac1p+\frac3{\ln2}$.
\end{proof}

Thus if $C_n\to\infty$, every coordinate of every extremal function has influence tending to zero. This is not true
of arbitrary functions with a large ratio $H/I$: the function $f(x)\,x_{n+1}$ has the same entropy as $f$, total
influence $I(f)+1$, and a coordinate with $D_{n+1}=0$.

\subsection*{The condition cannot be asked of every coordinate}

\begin{proposition}\label{prop:point}
Let $k$ be even, $g$ a bent function on $\{-1,1\}^k$, $y^*$ a point, and define $f$ on $\{-1,1\}\times\{-1,1\}^k$ by
$f(x_0,y)=g(y)$ for $y\neq y^*$ and $f(x_0,y^*)=x_0$. Then $\Inf_0(f)=2^{-k}$ and
\[
\frac{E_0(f)}{\Inf_0(f)}=k+\log_2e+O\!\left(2^{-k/2}\right).
\]
\end{proposition}

\begin{proof}
Since $f=g+\mathbf 1_{y^*}(y)\,(x_0-g(y^*))$ and $\mathbf 1_{y^*}=2^{-k}\sum_T\chi_T(y^*)\chi_T$, for
$T\subseteq[k]$ we get $b_T=2^{-k}\chi_T(y^*)$ and $a_T=\hat g(T)-2^{-k}g(y^*)\chi_T(y^*)$. Because $g$ is bent,
$\hat g(T)^2=2^{-k}$, so $b_T^2=4^{-k}$ and $a_T^2=2^{-k}(1\pm2^{-k/2})^2$. Hence $\Inf_0(f)=\sum_Tb_T^2=2^{-k}$, and
every $q_T=b_T^2/(a_T^2+b_T^2)$ equals $2^{-k}(1+O(2^{-k/2}))$, so $\log_2(1/q_T)=k+O(2^{-k/2})$. Using $h(q)=q\log_2\frac1q+q\log_2e+O(q^2)$,
\[
E_0(f)=\sum_T(a_T^2+b_T^2)\,h(q_T)=\sum_T b_T^2\Bigl(\log_2\frac1{q_T}+\log_2e+O(q_T)\Bigr)
=2^{-k}\bigl(k+\log_2e+O(2^{-k/2})\bigr).\qedhere
\]
\end{proof}

The ratio is unbounded while $I(f)\approx k/2$ and $H(f)\approx k$, so FEI is not in danger: the other
coordinates are good. For $k=2,4,6,8,10$ the exact values of $E_0/\Inf_0$ are $2.672$, $5.224$, $7.387$, $9.429$
and $11.439$. The same function shows that the ratio $J_0/\Inf_0$ tends to $2/\ln2$, which completes the proof of
Lemma~\ref{lem:jsd}.

\begin{remark}
A sufficient condition used in~\cite[Corollary 4.3]{GMP} is $\sum_T|\mu^+(T)-\mu^-(T)|\le K\Inf_i(f)$ for one
coordinate, which is Han's Question~3.1~\cite{Han}, since $|\mu^+(T)-\mu^-(T)|=4|a_Tb_T|$. This fails at every
coordinate of some read-once formulas, with ratio at least the depth~\cite{Shm}; there the one-coordinate entropy
step~\eqref{eq:step} still holds with $K=(5+\ln4)/\ln 2$. The erasure condition behaves differently: by Lemma~\ref{lem:sum}, every
nonconstant Boolean $f$ with $H(f)\le C\,I(f)$ has a coordinate with $E_i(f)\le C\,\Inf_i(f)$, so it cannot fail on
any class where FEI holds.
\end{remark}

\section{Computations}
The identity of Lemma~\ref{lem:split}, the bound of Lemma~\ref{lem:jsd} and the values in
Proposition~\ref{prop:point} were checked numerically: the identity on all Boolean functions of four variables and on
$40{,}000$ random functions each of five, six and seven variables at three densities, with $D_i$ computed from the
restricted truth tables and $E_i,J_i$ from the Fourier coefficients (largest discrepancy $7\cdot10^{-15}$); the
largest observed $J_i/\Inf_i$ was $2.806<2/\ln2\approx2.885$, the value of $J_0/\Inf_0$ for the function of
Proposition~\ref{prop:point} with $k=4$. On all functions of four variables
$\max_f\min_iE_i/\Inf_i=2.672$, and on the random sample of five variables it was $2.696$. The scripts are in the
repository accompanying this note.

\section*{Acknowledgements and funding}
This work was done by multiple AI agents, Claude (Anthropic) and Codex (OpenAI), connected to Math
Brain~\cite{MathBrainRepo} and working together inside the Brain framework~\cite{BrainFramework}, a shared, persistent
memory in which the agents record claims, proofs, failed attempts and methods, retrieve one another's results,
coordinate who works on what, and audit each other's work. For more on the details behind the Brain and Math Brain, see
the Math Brain paper~\cite{MathBrainPaper}, the Math Brain repository~\cite{MathBrainRepo} and the Brain framework
repository~\cite{BrainFramework}.

The research presented in this paper was supported by the European Research Council (ERC) under the European Union's
Horizon Europe research and innovation programme (grant agreement No.~101041711), by the Simons Foundation as part of
the Collaboration on the Mathematical and Scientific Foundations of Deep Learning, by Heights Labs, by Convex Nexus
Capital, by the Israel Science Foundation (grant number 2258/19), and by the Israel Science Foundation (ISF Grant
4101/25). Computational resources were provided in part by Google Cloud through its education credits programme (credit
530550103).

\section*{Declaration of competing interest}
The author holds equity in Heights Labs and in Convex Nexus Capital and has a management role at Convex Nexus Capital.
These entities are included in the funding declaration above.

\section*{Declaration of generative AI and AI-assisted technologies}
The paper comes out of Math Brain, an ongoing project on the use of AI systems for research mathematics, and its
results were developed with the Brain~\cite{BrainFramework} and Math Brain~\cite{MathBrainRepo}. In November 2026 a
Claude Fable agent (Anthropic), connected to the Brain and working alongside Codex agents running OpenAI's GPT-6.1 on
the same problem, derived the one-coordinate form of the conjecture, the averaging argument over coordinates and the
bent point-replacement example, and wrote the first numerical checks. Claude Opus 5.5 (Anthropic) searched the
literature, proved the sharp constant in Lemma~\ref{lem:jsd}, wrote independent code for every numerical claim,
prepared the manuscript and audited it. These checks are not peer review. The author directed the work and is
responsible for the content of the paper and for decisions concerning its dissemination.

\section*{Code and data availability}
\url{https://github.com/yspennstate/fourier-entropy-influence-one-coordinate}

The paper is part of the Math Brain project, whose public repository is
\url{https://github.com/yspennstate/math-brain-public}~\cite{MathBrainRepo}.

\appendix
\section{Notation}\label{app:notation}

Symbols are listed in the order in which they first appear, and the last column gives that place. The letters $g$ and
$p$ are reused, and each meaning has its own row.

\begin{longtable}{@{}>{\raggedright\arraybackslash}p{0.19\textwidth}>{\raggedright\arraybackslash}p{0.54\textwidth}>{\raggedright\arraybackslash}p{0.2\textwidth}@{}}
\toprule
Symbol & Meaning & First appears \\
\midrule
\endhead
$f$ & a Boolean function $\{-1,1\}^n\to\{-1,1\}$ & Section~\ref{sec:intro} \\
$\hat f(S)$, $\chi_S$ & the Fourier coefficient of $f$ at $S\subseteq[n]$, and the character $\prod_{i\in S}x_i$ & Section~\ref{sec:intro} \\
$\mu$ & the spectral distribution, $\mu(S)=\hat f(S)^2$ & Section~\ref{sec:intro} \\
$H(f)$ & the Fourier entropy $\sum_S\mu(S)\log_2(1/\mu(S))$ & Section~\ref{sec:intro} \\
$\Inf_i(f)$ & the influence of $i$, $\sum_{S\ni i}\mu(S)$ & Section~\ref{sec:intro} \\
$I(f)$ & the total influence $\sum_i\Inf_i(f)$ & Section~\ref{sec:intro} \\
$C$ & a universal constant, as in $H(f)\le C\,I(f)$ & Section~\ref{sec:intro} \\
$f^\pm$ & the restrictions $f_{i\to\pm1}$ & Section~\ref{sec:intro} \\
$D_i(f)$ & the entropy drop $H(f)-\tfrac12H(f^+)-\tfrac12H(f^-)$ & Section~\ref{sec:intro} \\
$K$ & a constant, as in~\eqref{eq:step} & Section~\ref{sec:intro} \\
$C_n$ & $\max H(f)/I(f)$ over nonconstant Boolean $f$ on $n$ variables & Section~\ref{sec:intro} \\
$E_i(f)$ & the erasure entropy $H(X_i\mid U\setminus\{i\})$, with $U$ and $X_i$ as below & Section~\ref{sec:intro} \\
$R_n$, $R'_n$ & $\max_f\min_iE_i(f)/\Inf_i(f)$ and $\max_f\min_iD_i(f)/\Inf_i(f)$, over the same functions as $C_n$ & Section~\ref{sec:intro} \\
$U$, $X_i$ & a random set with law $\mu$, and the indicator that $i\in U$ & Section~\ref{sec:split} \\
$a_T$, $b_T$ & $\hat f(T)$ and $\hat f(T\cup\{i\})$ for $T\subseteq[n]\setminus\{i\}$ & Section~\ref{sec:split} \\
$\mu^\pm$ & the spectral distributions of $f^\pm$, $\mu^\pm(T)=(a_T\pm b_T)^2$ & Section~\ref{sec:split} \\
$h(q)$ & the binary entropy $q\log_2\frac1q+(1-q)\log_2\frac1{1-q}$ & Section~\ref{sec:split} \\
$J_i(f)$, $\mathrm{JS}(p,q)$ & the Jensen--Shannon part $\mathrm{JS}(\mu^+,\mu^-)$ of $D_i(f)$, and the Jensen--Shannon divergence $H(\frac{p+q}2)-\frac12H(p)-\frac12H(q)$ of two laws $p$, $q$ & Section~\ref{sec:split} \\
$r$, $\theta$, $s$, $\sigma$ & the polar form $|a_T|=r\cos\theta$, $|b_T|=r\sin\theta$, $s=\sin2\theta$, and the sign of $a_Tb_T$, in the proof of Lemma~\ref{lem:jsd} & Lemma~\ref{lem:jsd} \\
$\varphi$, $F$ & $\varphi(s)=\frac12[(1+s)\ln(1+s)+(1-s)\ln(1-s)]$, and $F(\theta)=2\sin^2\theta-\varphi(\sin2\theta)$ & Lemma~\ref{lem:jsd} \\
$g$, $m$ & a Boolean function on $m\le n$ variables, in the induction & Section~\ref{sec:proof} \\
$p$ & the influence $\Inf_i(f)$ of the fixed coordinate &Proposition~\ref{prop:ext} \\
$k$, $g$, $y^*$, $x_0$ & an even number of variables, a bent function on $\{-1,1\}^k$, a point of $\{-1,1\}^k$, and the added coordinate & Proposition~\ref{prop:point} \\
$q_T$ & $b_T^2/(a_T^2+b_T^2)$ & Proposition~\ref{prop:point} \\
\bottomrule
\end{longtable}

\begin{thebibliography}{99}
\bibitem{ACKSW} S.~Arunachalam, S.~Chakraborty, M.~Kouck\'y, N.~Saurabh, R.~de~Wolf, Improved bounds on Fourier
entropy and min-entropy, ACM Trans. Comput. Theory 13 (2021), no.~4, 1--40, doi:10.1145/3470860; STACS 2020, arXiv:1809.09819.
\bibitem{Bes75} J.~Besag, Statistical analysis of non-lattice data, J. Roy. Statist. Soc. Ser. D (The Statistician)
24 (1975), 179--195.
\bibitem{CKLS} S.~Chakraborty, R.~Kulkarni, S.~V.~Lokam, N.~Saurabh, Upper bounds on Fourier entropy,
Electronic Colloquium on Computational Complexity TR13-052, 2013.
\bibitem{FK96} E.~Friedgut, G.~Kalai, Every monotone graph property has a sharp threshold, Proc. Amer. Math. Soc.
124 (1996), 2993--3002.
\bibitem{GMP} M.~J.~Gonz\'alez, P.~MacManus, M.~C.~Pereyra, Further evidence towards the Fourier
entropy-influence conjecture, arXiv:2606.00246, 2026.
\bibitem{Han} X.~Han, A new bound for the Fourier-entropy-influence conjecture, Combinatorica 45 (2025), no.~1,
Art.~4, doi:10.1007/s00493-024-00133-z; arXiv:2312.08271.
\bibitem{Han78} T.~S.~Han, Nonnegative entropy measures of multivariate symmetric correlations, Inform. Control 36
(1978), 133--156.
\bibitem{KKLMS} E.~Kelman, G.~Kindler, N.~Lifshitz, D.~Minzer, M.~Safra, Towards a proof of the
Fourier-entropy conjecture?, Geom. Funct. Anal. 30 (2020), 1097--1138.
\bibitem{KMS} N.~Keller, E.~Mossel, T.~Schlank, A note on the entropy/influence conjecture, Discrete Math. 312
(2012), 3364--3372.
\bibitem{LH25} P.~Li, G.~Han, Strengthening Han's Fourier entropy-influence inequality via an
information-theoretic proof, arXiv:2512.03117, 2025.
\bibitem{OT13} R.~O'Donnell, L.-Y.~Tan, A composition theorem for the Fourier entropy-influence conjecture,
ICALP 2013, Lecture Notes in Comput. Sci. 7965, 780--791, doi:10.1007/978-3-642-39206-1\_66; arXiv:1304.1347.
\bibitem{OWZ11} R.~O'Donnell, J.~Wright, Y.~Zhou, The Fourier entropy-influence conjecture for certain classes of
Boolean functions, ICALP 2011.
\bibitem{Pot26} V.~N.~Potapov, On the Fourier entropy-influence conjecture for Boolean plateaued functions,
arXiv:2609.21563, 2026.
\bibitem{Shm} Y.~Shmalo, Read-once formulas and a question of Han on the Fourier entropy-influence conjecture,
\url{https://github.com/yspennstate/han-question-read-once-formulas}, 2026.
\bibitem{ShmLTF} Y.~Shmalo, Fourier entropy of threshold functions and their compositions,
arXiv:2610.16069, 2026.
\bibitem{VW08} S.~Verd\'u, T.~Weissman, The information lost in erasures, IEEE Trans. Inform. Theory 54 (2008),
5030--5058.
\bibitem{WWW14} A.~Wan, J.~Wright, C.~Wu, Decision trees, protocols, and the Fourier entropy-influence conjecture,
ITCS 2014, arXiv:1312.3003.
\bibitem{MathBrainRepo} Y.~Shmalo, \emph{Math Brain: public snapshot, code and tests}, 2026,
\url{https://github.com/yspennstate/math-brain-public}.
\bibitem{MathBrainPaper} Y.~Shmalo, \emph{Math Brain: shared, auditable memory for multi-agent mathematical research},
October 2026, \url{https://github.com/yspennstate/math-brain-public/tree/main/papers/math-brain-2026-10-24}.
\bibitem{BrainFramework} Y.~Shmalo, \emph{Brain Framework}, public software repository, 2026,
\url{https://github.com/yspennstate/brain-framework-public}.
\end{thebibliography}

\end{document}
